\chapter{Validation et résultats}
\label{chap:validation}

Ce chapitre présente les interfaces réalisées, puis la stratégie de validation
(tests unitaires, d'intégration, d'API, RAG, \emph{streaming} et performance),
avant une analyse critique honnête des limites et un retour d'expérience. Les
emplacements de captures d'écran sont à renseigner manuellement à partir de
l'application en fonctionnement.

\section{Présentation des interfaces}
Pour chaque page, nous précisons l'objectif, les actions utilisateur et une
analyse fonctionnelle.

\subsection{Page d'accueil (Landing)}
\screenshotfig{Figure — Page d'accueil de DevPilot AI}{fig:cap-landing}
\textbf{Objectif.} Présenter la proposition de valeur et orienter vers la
connexion. \textbf{Actions.} Lecture des sections, navigation vers
\emph{Connexion}/\emph{Inscription}. \textbf{Analyse fonctionnelle.} Vitrine
toujours sombre, structurée en sections (problème, fonctionnalités, comparaison),
qui communique la différenciation (ancrage, traçabilité, observabilité).

\subsection{Connexion}
\screenshotfig{Figure — Interface de connexion}{fig:cap-login}
\textbf{Objectif.} Authentifier l'utilisateur. \textbf{Actions.} Saisie de
l'e-mail et du mot de passe, soumission. \textbf{Analyse fonctionnelle.} Émission
d'un jeton JWT, gestion des erreurs et redirection vers le tableau de bord ;
champs accessibles avec retours de validation.

\subsection{Inscription}
\screenshotfig{Figure — Interface d'inscription}{fig:cap-register}
\textbf{Objectif.} Créer un compte. \textbf{Actions.} Saisie du nom, de l'e-mail
et du mot de passe. \textbf{Analyse fonctionnelle.} Création d'un utilisateur de
rôle \texttt{user}, mot de passe haché, sans renvoi de données sensibles.

\subsection{Tableau de bord (personnel et administrateur)}
\screenshotfig{Figure — Tableau de bord personnel}{fig:cap-dashboard}
\textbf{Objectif.} Offrir un point d'entrée adapté au rôle. \textbf{Actions.}
Créer un espace, accéder aux documents et au chat ; pour l'administrateur,
consulter la vue d'exploitation. \textbf{Analyse fonctionnelle.} Le tableau de
bord est \emph{adapté au rôle} : les utilisateurs disposent d'un espace personnel
(création d'espaces, actions rapides), les administrateurs d'une vue d'opérations.

\subsection{Espaces de travail}
\screenshotfig{Figure — Gestion des espaces de travail}{fig:cap-workspaces}
\textbf{Objectif.} Organiser les connaissances par espace isolé. \textbf{Actions.}
Créer, ouvrir et gérer un espace et ses membres. \textbf{Analyse fonctionnelle.}
L'isolation par espace conditionne l'accès aux documents et conversations.

\subsection{Documents}
\screenshotfig{Figure — Gestion documentaire}{fig:cap-documents}
\textbf{Objectif.} Téléverser et suivre les documents. \textbf{Actions.}
Glisser-déposer, téléversement par lots, suivi de statut, relance.
\textbf{Analyse fonctionnelle.} Les statuts (\texttt{queued}, \texttt{processing},
\texttt{indexed}, \texttt{failed}) reflètent l'ingestion asynchrone ; les
documents en échec peuvent être relancés.

\subsection{Chat (streaming et citations)}
\screenshotfig{Figure — Interface de chat conversationnel}{fig:cap-chat}
\textbf{Objectif.} Poser des questions et recevoir des réponses citées.
\textbf{Actions.} Saisir une question, suivre la réponse en \emph{streaming},
consulter les sources. \textbf{Analyse fonctionnelle.} La réponse s'écrit jeton
par jeton ; les citations et sources sont reliées aux passages ; un accès à la
trace est offert aux administrateurs.

\subsection{Paramètres}
\screenshotfig{Figure — Paramètres}{fig:cap-settings}
\textbf{Objectif.} Configurer le compte et consulter l'état du système.
\textbf{Actions.} Consultation et ajustement des préférences.
\textbf{Analyse fonctionnelle.} Page personnelle accessible à tout utilisateur
authentifié.

\subsection{Tableau de bord d'administration (Operations Overview)}
\screenshotfig{Figure — Vue d'exploitation administrateur}{fig:cap-admin}
\textbf{Objectif.} Synthétiser la santé de la plateforme. \textbf{Actions.}
Consulter les métriques, surveiller les anneaux de santé et les listes de
surveillance. \textbf{Analyse fonctionnelle.} Indicateurs de santé, pipeline
documentaire, qualité RAG, latence des agents, usage et coût des modèles, et
listes d'actions/réponses/documents à risque.

\subsection{Tableau de bord Qualité}
\screenshotfig{Figure — Tableau de bord Qualité}{fig:cap-quality}
\textbf{Objectif.} Analyser la qualité des réponses. \textbf{Actions.} Lire les
distributions et les listes des pires réponses. \textbf{Analyse fonctionnelle.}
Carte de score, distributions de fidélité et d'hallucination, issues des réponses
et latence par agent.

\subsection{Tour de contrôle RAGOps}
\screenshotfig{Figure — Tour de contrôle RAGOps}{fig:cap-ragops}
\textbf{Objectif.} Superviser l'exploitation du pipeline. \textbf{Actions.}
Examiner la santé des services, la santé par espace et le centre des échecs.
\textbf{Analyse fonctionnelle.} Santé Qdrant/Redis/Celery, couverture
d'\emph{embedding}, file d'ingestion et relance des documents en échec.

\subsection{Observatoire des traces RAG}
\screenshotfig{Figure — Observatoire des traces (liste)}{fig:cap-traces}
\textbf{Objectif.} Explorer et filtrer les traces. \textbf{Actions.} Filtrer
(espace, utilisateur, fidélité, hallucination, dates), trier, ouvrir une trace.
\textbf{Analyse fonctionnelle.} Cartes de trace interactives avec score de santé,
distributions et bandeau « à investiguer ».

\subsection{Détail d'une trace (AI Trace Inspector)}
\screenshotfig{Figure — Inspecteur de traces}{fig:cap-trace-detail}
\textbf{Objectif.} Reconstruire le cycle de vie complet d'une requête.
\textbf{Actions.} Parcourir les neuf sections d'analyse. \textbf{Analyse
fonctionnelle.} Résumé, chronologie des agents, explorateurs de récupération, de
reclassement et de citations, centre d'évaluation, analyse du correcteur (\emph{diff}),
usage des modèles et graphe interactif.

\subsection{AI Execution Studio}
\screenshotfig{Figure — AI Execution Studio}{fig:cap-studio}
\textbf{Objectif.} Visualiser l'exécution du pipeline en temps réel.
\textbf{Actions.} Lancer une question, observer les dix étapes animées, rejouer.
\textbf{Analyse fonctionnelle.} Pipeline animé piloté par le flux SSE réel, avatar
réactif et contrôles de lecture (1×/2×/5×).

\subsection{Architecture Explorer}
\screenshotfig{Figure — Architecture Explorer}{fig:cap-archi}
\textbf{Objectif.} Explorer l'architecture de façon interactive. \textbf{Actions.}
Lancer une requête (flux animé), cliquer un n{\oe}ud pour l'inspecter.
\textbf{Analyse fonctionnelle.} Carte React Flow avec flux de requête animé,
inspecteur de n{\oe}ud (rôle, configuration réelle, métriques d'exécution) et
pages immersives par composant.

\section{Tests unitaires}
La couche de présentation est couverte par \textbf{39 tests} (Vitest et Testing
Library) portant sur les fonctions pures : algorithme de \emph{diff} mot à mot
(plus longue sous-séquence commune), score de santé de trace, mappage des tons,
construction des paramètres de filtre de l'observatoire, mise en \emph{buckets}
des distributions, et rendu des composants \texttt{StatusBadge}/\texttt{Badge}. Le
\emph{backend} dispose d'une suite \texttt{pytest}.

\screenshotfig{Figure — Exécution de la suite de tests unitaires}{fig:cap-tests}

\section{Tests d'intégration}
Les tests d'intégration valident la chaîne ingestion $\rightarrow$ indexation
$\rightarrow$ récupération sur un corpus contrôlé, en vérifiant que les fragments
téléversés deviennent récupérables avec des scores cohérents.

\section{Tests d'API}
Les contrats des \emph{endpoints} sont vérifiés (codes de statut, schémas,
autorisation) à l'aide de requêtes \texttt{curl} contre l'API réelle, en
particulier le rejet des accès non autorisés (401/403) et la non-divulgation des
secrets par \texttt{/system/ai-config}.

\section{Tests RAG}
La qualité est validée selon trois axes : récupération (pertinence des passages),
ancrage (réponse étayée par le contexte) et citation (références correctes).
L'agent évaluateur fournit une mesure continue de la fidélité, de la pertinence,
de la précision du contexte et du risque d'hallucination.

\section{Tests de streaming}
La diffusion SSE a été vérifiée par \texttt{curl} : les événements arrivent
\emph{progressivement} (\texttt{trace}, puis une succession de \texttt{token},
puis \texttt{evaluation}, \texttt{message}, \texttt{done}), confirmant un
\emph{streaming} réel de jetons après la correction du comportement de génération
simulée du fournisseur.

\section{Tests de performance}
Les routes pèsent de 110 à 175~kB environ ; les visualisations lourdes sont
chargées de façon différée ; les animations s'exécutent à 60~images/s grâce à
l'usage exclusif de \texttt{transform}/\texttt{opacity} ; le \emph{backend}
asynchrone et la recherche HNSW concourent à une expérience réactive. Le
tableau~\ref{tab:accept} synthétise les critères d'acceptation.

\begin{table}[H]
\centering\caption{Critères d'acceptation}\label{tab:accept}\small
\begin{tabularx}{\textwidth}{@{}X X c@{}}
\toprule
\textbf{Objet validé} & \textbf{Critère} & \textbf{Résultat} \\
\midrule
Compilation \emph{frontend} & \texttt{next build} sans erreur & \checkmark \\
Qualité du code & \texttt{eslint} sans avertissement & \checkmark \\
Tests unitaires & 39 tests au vert & \checkmark \\
Streaming & Jetons progressifs (SSE) & \checkmark \\
Réponse RAG & Ancrée et citée & \checkmark \\
Évaluation & Calculée et persistée & \checkmark \\
Traçabilité & Trace reconstructible & \checkmark \\
Sécurité & Aucune fuite de secret / inter-espaces & \checkmark \\
Tests e2e automatisés (Playwright) & Couverture des parcours & À venir \\
\bottomrule
\end{tabularx}
\end{table}

\section{Analyse critique}
Le système atteint un niveau de maturité notable, mais plusieurs points restent
perfectibles. La migration de la couche de données vers React Query est en cours
sur certaines pages historiques. L'évaluation repose sur des heuristiques et un
juge LLM plutôt que sur un cadre formel à large échelle. Le déploiement cible est
mono-machine. Enfin, la couverture de tests automatisés gagnerait à être étendue
aux parcours de bout en bout.

\section{Limitations du projet}
La comptabilité fine d'usage des modèles dépend des métadonnées exposées par
chaque fournisseur et n'est pas toujours disponible ; l'inspecteur affiche alors
honnêtement « non capturé ». Le multi-tenant reste au niveau de l'espace de
travail, sans cloisonnement organisationnel complet. Les visualisations de
« l'univers vectoriel » représentent des positions projetées et non les
coordonnées réelles en 768 dimensions, choix de visualisation assumé.

\section{Retour d'expérience}
Ce projet a confirmé que la valeur d'une plateforme d'IA d'entreprise ne réside
pas seulement dans la qualité des réponses, mais dans leur auditabilité et leur
exploitabilité. Concevoir pour l'observabilité dès l'origine, traiter les
exigences non fonctionnelles comme des exigences de premier ordre et raisonner
systématiquement en compromis ont été les enseignements majeurs, aux côtés de la
maîtrise opérationnelle d'une pile distribuée conteneurisée.

\section{Conclusion}
La stratégie de validation couvre l'ensemble des couches, du test unitaire à
l'acceptation, en accordant une attention particulière aux propriétés spécifiques
d'une plateforme d'IA : ancrage, citations, qualité mesurée, \emph{streaming} et
sécurité. Les résultats confirment que \devpilot{} satisfait ses objectifs, tout
en identifiant honnêtement ses axes d'amélioration.
